\chapter{Methodology}

\section{Prior-only statistical reference}

For symbol $i$ and observation $t$, volume is transformed as
\[
  x_{i,t}=\log(1+V_{i,t}),
\]
and the close-to-close return is
\[
  r_{i,t}=\frac{P_{i,t}-P_{i,t-1}}{P_{i,t-1}}.
\]
For either series $q\in\{x,r\}$, the reference mean and sample standard
deviation use the previous 60 valid observations:
\[
 \mu^q_{i,t}=\frac{1}{60}\sum_{j=1}^{60}q_{i,t-j},
 \qquad
 s^q_{i,t}=\sqrt{\frac{1}{59}\sum_{j=1}^{60}
 \left(q_{i,t-j}-\mu^q_{i,t}\right)^2}.
\]
The score is
\[
 z^q_{i,t}=\frac{q_{i,t}-\mu^q_{i,t}}
 {\max\left(s^q_{i,t},f_q\right)},
 \qquad f_x=0.10,\quad f_r=0.01.
\]

The observed value at $t$ is not part of its own mean or standard deviation.
The implementation achieves this through a one-observation shift. For volume,
the previous 60 log-volume observations are used. For return, missing values
created by the gap rule are removed before selecting the previous 60 valid
returns.

\section{Flag rules}

A volume anomaly requires $z^V_{i,t}>3.0$. The condition is one-sided because
the rule screens for unusually high volume. A price anomaly requires
$|z^R_{i,t}|>3.0$, so both unusually positive and unusually negative returns
can qualify. A combined anomaly requires both conditions on the same row. The
broad anomaly table is the union of the two individual rules.

\begin{table}[htbp]
\centering
\caption{Reference detector parameters}
\begin{tabularx}{0.94\textwidth}{@{}l l X@{}}
\toprule
Parameter & Value & Interpretation \\
\midrule
Reference length & 60 & Prior valid observations; observed row excluded \\
Volume threshold & $z^V>3.0$ & High-volume rule \\
Return threshold & $|z^R|>3.0$ & Two-sided return rule \\
Log-volume floor & 0.10 & Lower bound for the volume denominator \\
Return floor & 0.01 & Lower bound for the return denominator \\
Gap limit & 10 days & Longer-gap return omitted \\
OHLC policy & Raw & Source values retained and issues logged \\
\bottomrule
\end{tabularx}
\end{table}

The floors limit unstable ratios on nearly constant series. They do not make a
z-score a probability. The threshold and reference length are fixed in the
source. They were not adjusted using the event labels.

\section{Worked flag examples}

GIF on 22 October 2024 provides the event-window hit used in the evaluation.
The stored close is 0.44 TND, volume is 14,344, and the daily return is 10.00\%.
The volume z-score is 1.341 and the return z-score is 5.236. Therefore the row
is a price anomaly but not a volume anomaly. The same-symbol news search finds
no qualifying local match. The row remains outside the watchlist because the
watchlist requires a combined signal.

Carthage Cement (CC) on 30 October 2024 illustrates the combined route. The
stored close is 2.24 TND, volume is 3,343,436, and the return is 5.660\%. The
volume z-score is 4.069 and the return z-score is 4.558. Both signal fields are
true and no local news match is found. The row therefore enters the spike
watchlist route.

These examples also show why the union flag count and the watchlist count serve
different purposes. The union preserves every threshold crossing for review.
The watchlist selects only one narrow conjunction. A price-only event hit can
contribute to recall while remaining outside the watchlist.

\section{Sustained activity decline}

The decay detector compares total volume in a recent six-month interval with
the six months immediately before it. A symbol is eligible for the ratio rule
when baseline volume is at least 1,000 units. A decline case satisfies
\[
 \frac{\mathrm{recent\ volume}}{\mathrm{baseline\ volume}}<0.20.
\]
The detector also records symbols with at least 90 days since their last
positive-volume observation. The union of the decline and inactivity rules is
deduplicated by symbol.

The decay rule answers a different question from a daily z-score. It describes
a sustained change in trading activity rather than one extreme observation.
Its result is not added to the 1,414 daily union flags. Sixteen decay cases are
generated. Ten have a local news match and six do not. The six without a match
join the watchlist.

\section{News context and categories}

Daily flags use a same-symbol news interval of seven calendar days before or
after the flag. Decay cases use 60 days around the last active date. The
classifier keeps both the nearest item available on or before the reference
date and the nearest later item. This preserves the distinction between
information available at the time and retrospective context.

The status \code{NEWS_MATCH_FOUND} means that a qualifying stored item exists.
The status \code{NO_LOCAL_MATCH} means that none exists in the interval. The
matched-news categories are assigned through keyword rules. Their counts are
311 other news, 200 earnings results, 18 dividends, 13 governance items, four
capital increases, one buyout, and one suspension. The categories sum to 548,
the number of union flags with a news match.

\section{Watchlist and broad-market context}

The spike watchlist condition is
\[
 \mathrm{combined\ anomaly}\;\land\;\mathrm{no\ local\ news\ match}.
\]
Ten daily flags satisfy this condition. Six decay cases without a news match
are appended, so the direct and refined files each contain 16 rows.

Refinement adds a PX1 index return and score, other-symbol co-flag count, and
stock return minus index return. A dated row is tagged \code{MARKET_WIDE} when
the absolute index score is at least 2.0 or at least three other symbols are
flagged on the same date. Otherwise it is tagged \code{COMPANY_SPECIFIC}.
Decay rows have no single event date and retain a not-applicable tag.

The 16-row refined file contains four market-wide spike rows, six
company-specific spike rows, and six undated decay rows. These tags change
review order and context. They do not change the underlying detector result.

\section{Descriptive funnel}

The funnel assigns every union flag to one mutually exclusive outcome. A flag
can fail the combined-signal condition, fail the no-news condition, fail both,
or satisfy both admission conditions.

\begin{table}[htbp]
\centering
\caption{Daily-flag watchlist funnel}
\begin{tabularx}{0.88\textwidth}{@{}Xr@{}}
\toprule
Outcome & Flags \\
\midrule
Single signal only, without a news match & 856 \\
Combined signal, but with a news match & 11 \\
Single signal and a news match & 537 \\
Combined signal, without a news match & 10 \\
\midrule
Total & 1,414 \\
\bottomrule
\end{tabularx}
\end{table}

The funnel is descriptive. It measures the effect of the triage rule and does
not claim that excluded flags are false or that admitted rows are true events.
